
\section{1. Introduction}

\subsection{1.1 Motivation}

Transformer inference at long sequence length is dominated by the 
$\mathcal{O}(L^2)$ cost of softmax attention, but the actual \textbf{useful} 
content in a sequence is often much smaller than $L$. In classification 
tasks like SST-2, only a handful of tokens -- the sentiment-bearing words -- 
determine the output. In question answering, only tokens near the answer 
span are critical. The rest are redundant.

Prior work exploits this in three ways:

1. \textbf{Fixed heuristics.} Prune tokens by attention magnitude (Kim \& Cho, 
   2021), or by position (longformer-style). Simple but not task-aligned.

2. \textbf{Learned selectors.} Train a selector network on top of a frozen 
   transformer (Wang et al., 2021). Requires task-specific training.

3. \textbf{Full retraining.} Train the transformer with pruning in the loop 
   (Rao et al., 2021). Effective but expensive.

\textbf{The gap.} None of these is a \textbf{drop-in} solution that (a) learns which 
tokens to keep, (b) preserves pretrained weights exactly, and (c) works 
across flat, nested, and causal architectures. PTP fills this gap.

\subsection{1.2 Problem Statement}

Given a pretrained transformer with hidden dimension $d$ and $H$ heads, we 
seek a modified attention module $\mathcal{A}'$ that:

1. \textbf{Prunes tokens:} Retains only a fraction $\rho = k/L$ of tokens.
2. \textbf{Preserves [CLS]:} The classification token is always retained.
3. \textbf{Handles padding:} Invalid tokens are excluded.
4. \textbf{Preserves causality:} For decoders, no future information leaks.
5. \textbf{Drop-in:} Pretrained projections are copied exactly.
6. \textbf{Quality:} $\Delta\text{Acc} \leq 5\%$ on SST-2.

\subsection{1.3 Our Approach: PTP}

\textbf{Progressive Token Pruning} (PTP) introduces a **learnable scoring 
network**:

\[
s_i = \sigma\left(\text{MLP}(x_i)\right) \in (0, 1)
\tag{1.1}
\]

where $x_i$ is the hidden state of token $i$, and the MLP is:

\[
\text{MLP}(x) = W_2 \cdot \text{GELU}(W_1 x + b_1) + b_2
\tag{1.2}
\]

with $W_1 \in \mathbb{R}^{64 \times d}$, $W_2 \in \mathbb{R}^{1 \times 64}$.

\textbf{Top-$k$ selection.} Let $k = \lceil \rho L \rceil$. We keep the $k$ 
tokens with the highest scores:

\[
\mathcal{T} = \text{top-}k_i \, s_i
\tag{1.3}
\]

\textbf{Key design choices:}

1. \textbf{[CLS] preservation.} We set $s_0 = 10^4$ before top-$k$, guaranteeing 
   the first token is always retained.
2. \textbf{Padding masking.} Invalid tokens get $s_i = -10^4$ before top-$k$.
3. \textbf{Gather-based KV pruning.} Rather than masking, we gather the top-$k$ 
   KV tensors, reducing memory.
4. \textbf{Causal-aware attention.} For decoders, we mask queries that would 
   attend to future keys by comparing original positions.

\subsection{1.4 Contributions}

\begin{itemize}
\item \textbf{A learnable token-selection module.} PTP scores tokens with a small
\end{itemize}
  MLP and prunes them via top-$k$.
\begin{itemize}
\item \textbf{[CLS] preservation.} Guaranteed by score clamping, with no additional
\end{itemize}
  parameters or logic.
\begin{itemize}
\item \textbf{Causal-aware pruning.} We preserve causality by comparing original
\end{itemize}
  positions after pruning.
\begin{itemize}
\item \textbf{A drop-in wrapper.} Pretrained projections are copied exactly; only
\end{itemize}
  the score network and LoRA adapters are trainable.
\begin{itemize}
\item \textbf{A systematic evaluation.} Benchmarked on DistilBERT, RoBERTa, and
\end{itemize}
  GPT-2 on SST-2 with ablations over keep ratio and score network depth.
\begin{itemize}
\item \textbf{Token importance analysis.} We show that learned scores correlate with
\end{itemize}
  content-word importance and reveal systematic patterns.

\subsection{1.5 Paper Organization}

Section 2 reviews related work. Section 3 provides background on attention 
and token pruning. Section 4 presents PTP. Section 5 describes integration 
with four architectures. Section 6 details the experimental setup. Section 
7 reports results. Section 8 analyzes the learned token scores. Section 9 
discusses practical guidance. Section 10 lists limitations. Section 11 
concludes.



\subsection{Notation}
